\clearpage % Rozdziały zaczynamy od nowej strony.
\section{Dyskusja wyników}

TODO

% W niniejszym rozdziale przedstawiono przekrojową interpretację wyników zaprezentowanych w rozdziale 5. Celem dyskusji jest powiązanie odpowiedzi na poszczególne pytania badawcze, wskazanie zależności między obserwowanymi zjawiskami oraz ocena znaczenia praktycznego uzyskanych rezultatów dla zastosowania agentów LLM w zadaniach DevOps.

% \subsection{Synteza odpowiedzi na pytania badawcze}

% % TODO: Zsyntetyzować odpowiedzi na RQ1--RQ5 w jednej narracji.
% % Proponowany zakres:
% % - RQ1: realna skuteczność end-to-end i jej granice,
% % - RQ2: wpływ złożoności architektury na spadek skuteczności,
% % - RQ3: relacja między jakością artefaktów a instrukcjami jakościowymi,
% % - RQ4: znaczenie zmienności procesu mimo deterministycznych ustawień,
% % - RQ5: podatność na manipulację jako ograniczenie zaufania do wyników.

% \subsection{Porównanie modeli}

% % TODO: Przeprowadzić analizę przekrojową modeli między RQ1--RQ5.
% % Proponowany zakres:
% % - porównanie skuteczności funkcjonalnej między modelami,
% % - różnice w jakości generowanych artefaktów,
% % - różnice w powtarzalności procesu,
% % - różnice w odporności na manipulację kontekstem,
% % - kompromisy między skutecznością, jakością i bezpieczeństwem.

% \subsection{Interpretacja ograniczeń modeli LLM}

% % TODO: Zinterpretować główne źródła błędów i ograniczeń modeli.
% % Proponowany zakres:
% % - ograniczone rozumienie kontekstu runtime i środowiska wykonawczego,
% % - problemy z zależnościami między usługami i konfiguracją infrastrukturalną,
% % - halucynacje i nadmierne uogólnianie jako mechanizm błędów,
% % - zależność jakości wyniku od kompletności i spójności sygnałów wejściowych.

% \subsection{Rola prompt engineeringu}

% % TODO: Wyjaśnić, dlaczego rozszerzony prompt poprawia jakość wyników.
% % Proponowany zakres:
% % - redukcja przestrzeni dopuszczalnych rozwiązań,
% % - wymuszenie uwzględniania dobrych praktyk,
% % - wpływ na jakość artefaktów i spójność decyzji,
% % - praktyczne implikacje dla projektowania instrukcji systemowych.

% \subsection{Niezawodność i deterministyczność procesu agentowego}

% % TODO: Omówić znaczenie wyników RQ4 w szerszym kontekście.
% % Proponowany zakres:
% % - dlaczego temperature=0 nie gwarantuje pełnej deterministyczności,
% % - wpływ architektury agentowej i narzędzi na zmienność wyników,
% % - różnica między pojedynczym wywołaniem modelu a wieloetapowym procesem agentowym,
% % - konsekwencje praktyczne dla automatyzacji i powtarzalności pipeline'ów.

% \subsection{Aspekty bezpieczeństwa i podatność na manipulację}

% % TODO: Omówić implikacje wyników RQ5 dla bezpieczeństwa praktycznych wdrożeń.
% % Proponowany zakres:
% % - znaczenie podatności na mylący kontekst repozytorium,
% % - repo poisoning i prompt injection jako realne ryzyka,
% % - wpływ na zaufanie do wygenerowanych artefaktów,
% % - potrzeba mechanizmów obronnych i dodatkowej weryfikacji.

% \subsection{Porównanie z literaturą}

% % TODO: Odnieść uzyskane wyniki do literatury omówionej w rozdziale 2.
% % Proponowany zakres:
% % - które obserwacje potwierdzają wcześniejsze prace,
% % - gdzie wyniki wskazują silniejsze ograniczenia praktyczne,
% % - różnice wynikające z zastosowania architektury agentowej i walidacji end-to-end.

% \subsection{Znaczenie praktyczne dla automatyzacji DevOps}

% % TODO: Zebrać najważniejsze implikacje praktyczne w jednym miejscu.
% % Proponowany zakres:
% % - zastosowania w prostszych projektach i jako wsparcie inżyniera,
% % - ograniczenia przy aplikacjach wielousługowych,
% % - potrzeba walidacji statycznej i testów runtime,
% % - warunki bezpiecznego użycia w środowisku organizacyjnym.

% \subsection{Ograniczenia badania i zagrożenia dla trafności}

% % TODO: Opisać ograniczenia metodologiczne i ich wpływ na interpretację wyników.
% % Proponowany zakres:
% % - dobór repozytoriów i liczba powtórzeń,
% % - ograniczony zestaw modeli,
% % - środowisko wykonawcze i pipeline walidacyjny,
% % - brak walidacji runtime w RQ3 i RQ5,
% % - możliwa zależność wyników od promptu i użytych narzędzi agenta.
